% One-page summary panel: abstract, keywords and the proposal at a glance.
\thispagestyle{scrheadings}
\addchap*{Summary}
\begin{summarypanel}
{\sffamily\small\color{accent}\textsc{in one sentence}\par}
\vspace{.4em}
{\large\itshape This thesis asks how a city can learn, from a few hundred
cheap sensors, where to put the next one so that street-level heat is
mapped well enough to protect the people most exposed to it.\par}

\vspace{1.2em}
Heat is the deadliest weather hazard in European cities, yet official
temperature records come from a handful of airport and rooftop stations
that miss the street canyons, courtyards and bus stops where people
actually spend hot afternoons. Low-cost sensor networks close part of
this gap, but they are placed by convenience rather than by need: on
lampposts that happen to have power, in neighbourhoods that volunteer.

The proposed research treats sensor placement as a sequential learning
problem. A Gaussian-process model of the urban temperature field, built
on land-cover and morphology covariates, scores each candidate site by
how much it would reduce uncertainty where vulnerable residents live.
The thesis develops this exposure-weighted acquisition rule
(\ref{rq:acq}), tests whether models trained in one city transfer to
another (\ref{rq:transfer}), and measures the gain in a live
deployment of 60 mobile units in Aldmere (\ref{rq:field}).

Preliminary work on 214 existing sensors shows that exposure-weighted
placement reaches the mapping error of random placement with roughly
40\,\% fewer sensors.

\vspace{1em}
{\color{accent!50}\rule{\linewidth}{.3pt}}
\vspace{.4em}

\begin{tabularx}{\linewidth}{@{}>{\raggedright\arraybackslash}p{30mm}X@{}}
\fieldlabel{Keywords} & urban heat, Gaussian processes, active learning,
  sensor networks, environmental justice \\
\fieldlabel{Questions} & three research questions, Chapter~\ref{ch:questions} \\
\fieldlabel{Duration} & four years, February 2025 to January 2029 \\
\fieldlabel{Funding} & Aldmere Doctoral Fellowship; field units from
  the City of Aldmere Climate Office \\
\fieldlabel{Data} & 214 municipal sensors (2021--2025), 3 m land cover,
  LiDAR building heights, census vulnerability index \\
\end{tabularx}
\end{summarypanel}

% At a glance: key numbers, filling the page below the panel.
\vspace{1em}
\noindent{\sffamily\small\color{accent}\textsc{at a glance}\par}
\vspace{.1em}
\noindent{\color{accent!50}\rule{\linewidth}{.3pt}}
\vspace{.5em}

\newcommand*{\glancefigure}[2]{%
  {\sffamily\Large\color{accent}#1\par}%
  {\footnotesize\raggedright #2\par}}
\noindent{\renewcommand{\arraystretch}{1}%
\begin{tabularx}{\linewidth}{@{}*{4}{>{\raggedright\arraybackslash}X}@{}}
\glancefigure{214}{municipal sensors behind the preliminary model} &
\glancefigure{0.82\,K}{leave-one-out RMSE on summer afternoons} &
\glancefigure{40\,\%}{fewer sensors for the error of random placement} &
\glancefigure{60}{mobile units in the 2027 field campaign} \\
\end{tabularx}\par}
\vfill
\clearpage
